\chapter{प्रथमः प्रपाठकः}

% --- Kandika 1 ---
\dnsub{॥ प्रथमा कण्डिका ॥}
\resetShloka % Resets counter to 1 for new Kandika

\twolineshloka{अथातः स्वरशास्त्राणां सर्वेषां वेदनिश्चयम्}
{उच्चनीचविशेषाद्धि स्वरान्यत्वं प्रवर्तते}

\twolineshloka{आर्चिकं गाथिकं चैव सामिकं च स्वरान्तरम्}
{कृतान्ते स्वरशास्त्राणां प्रयोक्तव्यं विशेषतः}

\twolineshloka{एकान्तरः स्वरो ह्यृक्षु गाथासु द्व्यन्तरः स्वरः}
{सामसु त्र्यन्तरं विद्यादेतावत् स्वरतोऽन्तरम्}

\twolineshloka{ऋक्सामयजुरङ्गानि ये यज्ञेषु प्रयुञ्जते}
{अविज्ञानाद्धि शास्त्राणां तेषां भवति विस्वरः}

\twolineshloka{मन्त्रो हीनः स्वरतो वर्णतो वा मिथ्याप्रयुक्तो न तमर्थमाह}
{स वाग्वज्रो यजमानं हिनस्ति यथेन्द्रशत्रुः स्वरतोऽपराधात्}

\twolineshloka{प्रहीणः स्वरवर्णाभ्यां यो विमन्त्रः प्रयुज्यते}
{यज्ञेषु यजमानस्य दहत्यायुः प्रजां पशून्}

\twolineshloka{उरः कण्ठः शिरश्चेति स्थानानि त्रीणि वाङ्मये}
{सवनान्याहुरेतानि साम वाप्यर्थतोऽन्तरम्}

\twolineshloka{उरः सप्तविचारं स्यात् तथा कण्ठस्तथा शिरः}
{न च सप्तोरसि व्यक्तास्तथा प्रावचनो विधिः}

\twolineshloka{कठकालापवृत्तेषु तैत्तिरीयाध्वरकेषु च}
{ऋग्वेदे सामवेदे च वक्तव्यः प्रथमः स्वरः}

\twolineshloka{ऋग्वेदस्तु द्वितीयेन तृतीयेन च वर्तते}
{उच्चमध्यमसङ्घातस्वरो भवति पार्थिवः}

\twolineshloka{तृतीयप्रथमक्रुष्टान् कुर्वन्त्याध्वरकाः स्वरान्}
{द्वितीयाद्यांस्तु मन्द्रान्तांस्तैत्तिरीयाश्चतुरः स्वरान्}

\twolineshloka{प्रथमश्च द्वितीयश्च तृतीयोऽथ चतुर्थकः}
{मन्द्रः क्रुष्टो ह्यतिस्वार एतान् कुर्वन्ति सामगाः}

\twolineshloka{द्वितीयप्रथमावेतौ ताण्डिभाल्लविनां स्वरौ}
{तथा शतपथावेतौ स्वरौ वाजसनेयिनाम्}

\twolineshloka{एते विशेषतः प्रोक्ता स्वरा वै सार्ववैदिकाः}
{इत्येतच्चरितं सर्वं स्वराणां सार्ववैदिकम्}

\closesub

% --- Kandika 2 ---
\dnsub{॥ द्वितीया कण्डिका ॥}
\resetShloka

\twolineshloka{सामवेदे तु वक्ष्यामि स्वराणां चरितं यथा}
{अल्पग्रन्थं प्रभूतार्थं श्रव्यं वेदाङ्गमुत्तमम्}

\twolineshloka{तानरागस्वरग्राममूर्छनानां तु लक्षणम्}
{पवित्रं पावनं पुण्यं नारदेन प्रकीर्तितम्}

\twolineshloka{शिक्षामाहुर्द्विजातीनामृग्यजुःसामलक्षणम्}
{नारदीयमशेषेण निरुक्तमनुपूर्वशः}

\twolineshloka{सप्त स्वरास्त्रयो ग्रामा मूर्छनास्त्वेकविंशतिः}
{ताना एकोनपञ्चाशदित्येतत् स्वरमण्डलम्}

\twolineshloka{षड्जश्च ऋषभश्चैव गान्धारो मध्यमस्तथा}
{पञ्चमो धैवतश्चैव निषादः सप्तमः स्वरः}

\twolineshloka{षड्जमध्यमगान्धारास्त्रयो ग्रामाः प्रकीर्तिताः}
{भूर्लोकाज्जायते षड्जो भुवर्लोकाच्च मध्यमः}

\twolineshloka{स्वर्गान्नान्यत्र गान्धारो नारदस्य मतं यथा}
{स्वररागविशेषेण ग्रामरागा इति स्मृताः}

\twolineshloka{विंशतिं मध्यमग्रामे षड्जग्रामे चतुर्दश}
{तानान् पञ्चदशेच्छन्ति गान्धारग्राममाश्रितान्}

\twolineshloka{नन्दी विशाला सुमुखी चित्रा चित्रवती सुखा}
{बला या चाथ विज्ञेया देवानां सप्त मूर्छनाः}

\twolineshloka{आप्यायनी विश्वभृता चन्द्रा हेमा कपर्दिनी}
{मैत्री च बार्हती चैव पितृणां सप्त मूर्छनाः}

\twolineshloka{षड्जे तूत्तरमन्द्रा स्यादृषभे चाभिरुद्गता}
{अश्वक्रान्ता तु गान्धारे तृतीया मूर्छना स्मृता}

\twolineshloka{मध्यमे खलु सौवीरा हृष्यका पञ्चमे स्वरे}
{धैवते चापि विज्ञेया मूर्छना तूत्तरायता}

\twolineshloka{निषादे रजनी विद्याद् ऋषीणां सप्त मूर्छनाः}
{उपजीवन्ति गन्धर्वा देवानां सप्त मूर्छनाः}

\twolineshloka{पितृणां मूर्छनाः सप्त तथा यक्षा न संशयः}
{ऋषीणां मूर्छनाः सप्त यास्त्विमा लौकिकाः स्मृताः}

\twolineshloka{षड्जः प्रीणाति वै देवानृषीन् प्रीणाति चर्षभः}
{पितृन् प्रीणाति गान्धारो गन्धर्वान् मध्यमः स्वरः}

\twolineshloka{देवान् पितृनृषींश्चैव स्वरः प्रीणाति पञ्चमः}
{यक्षान् निषादः प्रीणाति भूतग्रामं च धैवतः}

\closesub

% --- Kandika 3 (Prose mix) ---
\dnsub{॥ तृतीया कण्डिका ॥}
\resetShloka

\onelineshloka{गानस्य तु दशविधा गुणवृत्तिः}

% Prose section - using \gadyam command
\gadyam{तद्यथा- रक्तं पूर्णमलङ्कृतं प्रसन्नं व्यक्तं विक्रुष्टं श्लक्ष्णं समं सुकुमारं मधुरमिति गुणाः \nextShloka{\ddanda}}

\gadyam{तत्र \textbf{रक्तं} नाम- वेणुवीणास्वराणामेकीभावे रक्तमित्युच्यते \nextShloka{\ddanda}}

\gadyam{\textbf{पूर्णं} नाम- स्वरश्रुतिपूरणाच्छन्दःपादाक्षरसंयोगात् पूर्णमित्युच्यते \nextShloka{\ddanda}}

\gadyam{\textbf{अलङ्कृतं} नाम- उरसि शिरसि कण्ठयुक्तमित्यलङ्कृतमित्युच्यते \nextShloka{\ddanda}}

\gadyam{\textbf{प्रसन्नं} नाम- अपगतगद्गदं निर्विशङ्कं प्रसन्नमित्युच्यते \nextShloka{\ddanda}}

% This is the long line causing overflow. \gadyam will fix it automatically.
\gadyam{\textbf{व्यक्तं} नाम- पदपदार्थप्रकृतिविकारागमलोपकृत्तद्धितसमासधातुनिपातोपसर्गस्वरलिङ्गवृत्तिवार्त्तिकविभक्त्यर्थवचनानां सम्यगुपपादने व्यक्तमित्युच्यते \nextShloka{\ddanda}}

\gadyam{\textbf{विक्रुष्टं} नाम- उच्चैरुच्चारितं व्यक्तपदाक्षरं विक्रुष्टमित्युच्यते \nextShloka{\ddanda}}

\gadyam{\textbf{श्लक्ष्णं} नाम- अद्रुतमविलम्बितमुच्चनीचप्लुतसमाहारं लयतालोपनयादिभिरुपपादनैः श्लक्ष्णमित्युच्यते \nextShloka{\ddanda}}

\gadyam{\textbf{समं} नाम- आवापनिर्वापप्रदेशं प्रति अन्तरस्थानानां समासः सममित्युच्यते \nextShloka{\ddanda}}

\gadyam{\textbf{सुकुमारं} नाम- मृदुपदं वर्णस्वरकुहरणयुक्तं सुकुमारमित्युच्यते \nextShloka{\ddanda}}

\gadyam{\textbf{मधुरं} नाम- स्वभावोपनीतललितपदाक्षरगुणसमृद्धं मधुरमित्युच्यते \nextShloka{\ddanda}}

\onelineshloka{एवमेतैर्दशभिर्गुणैर्युक्तं गानं भवति}

\onelineshloka{एतद्विपरीता गीतिदोषा उच्यन्ते}

भवन्ति चात्र श्लोकाः-

\twolineshloka{शङ्कितं भीतमुद्घुष्टमव्यक्तमनुनासिकम्}
{काकस्वरं शिरसि गतं तथा स्थानविवर्जितम्}

\twolineshloka{विस्वरं विरसं चैव विश्लिष्टं विषमाहतम्}
{व्याकुलं तालहीनं च गीतिदोषाश्चतुर्दश}

\twolineshloka{आचार्याः सममिच्छन्ति पदच्छेदं तु पण्डिताः}
{स्त्रियो मधुरमिच्छन्ति विक्रुष्टमितरे जनाः}

\closesub

% --- Kandika 4 ---
\dnsub{॥ चतुर्थी कण्डिका ॥}
\resetShloka

\twolineshloka{पद्मपत्रप्रभः षड्ज ऋषभः शुकपिञ्जरः}
{कनकाभस्तु गान्धारो मध्यमः कुन्दसन्निभः}

\twolineshloka{पञ्चमस्तु भवेत् कृष्णः पीतकं धैवतं विदुः}
{निषादः सर्ववर्णः स्यादित्येताः स्वरवर्णताः}

\twolineshloka{पञ्चमो मध्यमः षड्ज इत्येते ब्राह्मणाः स्मृताः}
{ऋषभो धैवतश्चापीत्येतौ हि क्षत्रियावुभौ}

\twolineshloka{गान्धारश्च निषादश्च वैश्यावेव हि स्मृतौ}
{शूद्रत्वं विद्धि चार्धेन पतितत्वान् न संशयः}

\twolineshloka{ऋषभोत्थितषड्जहतो धैवतसहितश्च पञ्चमो यत्र}
{निपतति मध्यमरागे तं निषादं षाडवं विद्यात्}

\twolineshloka{यदि पञ्चमो विरमते गान्धारश्चान्तरस्वरो भवति}
{ऋषभो निषादसहितस्तं पञ्चममीदृशं विद्यात्}

\twolineshloka{गान्धारस्याधिपत्येन निषादस्य गतागतैः}
{धैवतस्य च दौर्बल्यान् मध्यमग्राम उच्यते}

\twolineshloka{ईषत् स्पृष्टो निषादस्तु गान्धारश्चाधिको भवेत्}
{धैवतः कम्पितो यत्र षड्जग्रामं तु निर्दिशेत्}

\twolineshloka{अन्तरस्वरसंयुक्ता काकलिर्यत्र दृश्यते}
{तं तु साधारितं विद्यात् पञ्चमस्थं तु कैशिकम्}

\twolineshloka{कैशिकं भावयित्वा तु स्वरैः सर्वैः समन्ततः}
{यस्मात्तु मध्यमे न्यासस्तस्मात् कैशिकमध्यमः}

\twolineshloka{काकलिर्दृश्यते यत्र प्राधान्यं पञ्चमस्य तु}
{कश्यपः कैशिकं प्राह मध्यमग्रामसम्भवम्}

\twolineshloka{गेति गेयं विदुः प्राज्ञा धेति कारुप्रवादनम्}
{वेति वाद्यस्य संज्ञानं गान्धर्वस्य विरोचनम्}

\closesub

% --- Kandika 5 ---
\dnsub{॥ पञ्चमी कण्डिका ॥}
\resetShloka

\twolineshloka{यः सामगानां प्रथमः स वेणोर्मध्यमः स्वरः}
{यो द्वितीयः स गान्धारस्तृतीयस्त्वृषभः स्मृतः}

\twolineshloka{चतुर्थः षड्ज इत्याहुः पञ्चमो धैवतो भवेत्}
{षष्ठो निषादो विज्ञेयः सप्तमः पञ्चमः स्मृतः}

\twolineshloka{षड्जं वदति मयूरो गावो रम्भन्ति चर्षभम्}
{अजाविके तु गान्धारं क्रौञ्चो वदति मध्यमम्}

\twolineshloka{पुष्पसाधारणे काले कोकिला वक्ति पञ्चमम्}
{अश्वस्तु धैवतं वक्ति निषादं वक्ति कुञ्जरः}

\twolineshloka{कण्ठादुत्तिष्ठते षड्जः शिरसस्त्वृषभः स्मृतः}
{गान्धारस्त्वनुनासिक्य उरसो मध्यमः स्वरः}

\twolineshloka{उरसः शिरसः कण्ठादुत्थितः पञ्चमः स्वरः}
{ललाटाद्धैवतं विद्यान्निषादं सर्वसन्धिजम्}

\twolineshloka{नासां कण्ठमुरस्तालु जिह्वां दन्तांश्च संश्रितः}
{षड्भिः सञ्जायते यस्मात्तस्मात् षड्ज इति स्मृतः}

\twolineshloka{वायुः समुत्थितो नाभेः कण्ठशीर्षसमाहतः}
{नर्दत्यूषभवद्यस्मात्तस्मादृषभ उच्यते}

\twolineshloka{वायुः समुत्थितो नाभेः कण्ठशीर्षसमाहतः}
{नासागन्धावहः पुण्यो गान्धारस्तेन हेतुना}

\twolineshloka{वायुः समुच्छ्रितो नाभेरुरोहृदि समाहतः}
{नाभिं प्राप्तो महानादो मध्यमत्वं समश्नुते}

\twolineshloka{वायुः समुत्थितो नाभेरुरोहृत्कण्ठशिरोहतः}
{पञ्चस्थानोत्थितस्यास्य पञ्चमत्वं विधीयते}

\twolineshloka{धैवतं च निषादं च वर्जयित्वा स्वरद्वयम्}
{शेषांस्त्वन्यान् पञ्च स्वराँस्त्वन्यान् पञ्चस्थानोच्छ्रितान् विदुः}

\twolineshloka{अग्निगीतः स्वरः षड्ज ऋषभो ब्रह्मणोच्यते}
{सोमेन गीतो गान्धारो विष्णुना मध्यमः स्वरः}

\twolineshloka{पञ्चमस्तु स्वरो गीतो नारदेन महात्मना}
{धैवतश्च निषादश्च गीतौ तुम्बुरुणा स्वरौ}

\twolineshloka{आद्यस्य दैवतं ब्रह्मा षड्जस्याप्युच्यते बुधैः}
{तीक्ष्णदीप्तिप्रकाशत्वादृषभस्य हुताशनः}

\twolineshloka{गावः प्रगीते तुष्यन्ति गान्धारस्तेन हेतुना}
{श्रुत्वा चैवोपतिष्ठन्ति सौरभेया न संशयः}

\twolineshloka{सोमस्तु पञ्चमस्यापि दैवतं ब्रह्मराट् स्मृतम्}
{निर्हासो यच्च वृद्धिश्च ग्राममासाद्य सोमवत्}

\twolineshloka{अतिसन्धीयते यस्मादेतान् पूर्वोत्थितान् स्वरान्}
{तस्मादस्य स्वरस्यापि धैवतत्वं विधीयते}

\twolineshloka{निषीदन्ति स्वरा यस्मान्निषादस्तेन हेतुना}
{सर्वांश्चाभिभवत्येष यदादित्योऽस्य दैवतम्}

\closesub

% --- Kandika 6 ---
\dnsub{॥ षष्ठी कण्डिका ॥}
\resetShloka

\twolineshloka{दारवी गात्रवीणा च द्वे वीणे गानजातिषु}
{सामिकी गात्रवीणा तु तस्याः शृणुत लक्षणम्}

\twolineshloka{गात्रवीणा तु सा प्रोक्ता यस्यां गायन्ति सामगाः}
{स्वरव्यञ्जनसंयुक्ता अङ्गुल्यङ्गुष्ठरञ्जिता}

\twolineshloka{हस्तौ सुसंयतौ धार्यौ जानुभ्यामुपरि स्थितौ}
{गुरोरनुकृतिं कुर्याद्यथा ज्ञानमतिर्भवेत्}

\twolineshloka{प्रणवं प्राक् प्रयुञ्जीत व्याहृतीस्तदनन्तरम्}
{सावित्रीं चानुवचनं ततो वृत्तान्तमारभेत्}

\twolineshloka{प्रसार्य चाङ्गुलीः सर्वा रोपयेत् स्वरमण्डलम्}
{न चाङ्गुलीभिरङ्गुष्ठमङ्गुष्ठेनाङ्गुलीः स्पृशेत्}

\twolineshloka{विरला नाङ्गुलीः कुर्यान्मूले चैना न संस्पृशेत्}
{अङ्गुष्ठाग्रेण ता नित्यं मध्यमे पर्वणि स्पृशेत्}

\twolineshloka{मात्रद्विमात्रवृद्धानां विभागार्थं विभागवित्}
{अङ्गुलीभिर्द्विमात्रं तु पाणेः सव्यस्य दर्शयेत्}

\twolineshloka{त्रिरेखा यत्र दृश्येत सन्धिं तत्र विनिर्दिशेत्}
{स पर्व इति विज्ञेयः शेषमन्तरमन्तरम्}

\twolineshloka{यवान्तरं तु सामसु ऋक्षु कुर्यात् तिलान्तरम्}
{स्वरान् मध्यमपर्वसु सुनिविष्टान् निवेशयेत्}

\twolineshloka{न चात्र कम्पयेत् किञ्चिदङ्गस्यावयवं बुधः}
{अधस्तनं मृदु न्यस्य हस्तमास्ते यथाक्रमम्}

\twolineshloka{कूर्मोऽङ्गानीव संहृत्य चेष्टां दृष्टिं दिशं मनः}
{स्वस्थः प्रशान्तो निर्भीतो वर्णानुच्चारयेद् बुधः}

\twolineshloka{नासिकायास्तु पूर्वेण हस्तं गोकर्णवद्धरेत्}
{निवेश्य दृष्टिं हस्ताग्रे शास्त्रार्थमनुचिन्तयेत्}

\twolineshloka{सममुच्चारयेद्वाक्यं हस्तेन च मुखेन च}
{यथैवोच्चारयेद्वर्णांस्तथैवैनान् समापयेत्}

\twolineshloka{नात्याहन्यान्न निर्हण्यान्न प्रगायेन्न कम्पयेत्}
{समं सामानि गायेत व्योम्नि श्येनगतिर्यथा}

\twolineshloka{यथाप्सु चरतां मार्गो मीनानां नोपलभ्यते}
{आकाशे वा विहङ्गानां तद्वत् स्वरगता श्रुतिः}

\twolineshloka{यथा दधनि सर्पिः स्यात् काष्ठस्थो वा यथानलः}
{प्रयत्नेनोपलभ्येत तद्वत् स्वरगता श्रुतिः}

\twolineshloka{स्वरात् स्वरं सङ्क्रामंस्तु स्वरसन्धिमनुल्बणम्}
{अविच्छिन्नं समं कुर्यात् सूक्ष्मं छायाऽऽतपोपमम्}

\twolineshloka{अनागतमतिक्रान्तं विच्छिन्नं विषमाहतम्}
{तन्वन्तमस्थितान्तं च वर्जयेत् कर्षणं बुधः}

\twolineshloka{स्वरः स्थानाच्च्युतो यस्तु स्वस्थानमतिवर्तते}
{विस्वरं सामगा ब्रूयुर्विरक्तमिति वीणिनः}

\twolineshloka{अभ्यासार्थे द्रुतां वृत्तिं प्रयोगार्थे तु मध्यमाम्}
{शिष्याणामुपदेशार्थे कुर्याद् वृत्तिं विलम्बिताम्}

\twolineshloka{गृहीतग्रन्थ एवं तु ग्रन्थोच्चारणशिक्षकान्}
{हस्तेनाध्यापयेच्छिष्यान् शैक्षेण विधिना द्विजः}

\closesub

% --- Kandika 7 ---
\dnsub{॥ सप्तमी कण्डिका ॥}
\resetShloka

\twolineshloka{क्रुष्टस्य मूर्धनि स्थानं ललाटे प्रथमस्य तु}
{भ्रुवोर् मध्ये द्वितीयस्य तृतीयस्य तु कर्णयोः}

\twolineshloka{कण्ठः स्थानं तु तुर्यस्य मन्द्रस्योरसि तूच्यते}
{अतिस्वारस्य नीचस्य हृदि स्थानं विधीयते}

\twolineshloka{अङ्गुष्ठस्योत्तमे क्रुष्टोऽङ्गुष्ठे तु प्रथमः स्वरः}
{प्रदेशिन्यां तु गान्धार ऋषभस्तदनन्तरम्}

\twolineshloka{अनामिकायां षड्जस्तु कनिष्ठायां च धैवतः}
{तस्याधस्ताच्च यान्या तु निषादं तत्र विन्यसेत्}

\twolineshloka{अपर्वत्वादसंज्ञत्वादव्ययत्वाच्च नित्यशः}
{मन्द्रो हिनुहिभूतस्तु परिस्वार इति स्मृतः}

\twolineshloka{क्रुष्टेन देवा जीवन्ति प्रथमेन तु मानुषाः}
{पशवस्तु द्वितीयेन गन्धर्वाप्सरसस्त्वनु}

\twolineshloka{आण्डजाः पितरश्चैव चतुर्थस्वरजीविनः}
{मन्द्रं चैवोपजीवन्ति पिशाचासुरराक्षसाः}

\twolineshloka{अतिस्वारेण नीचेन जगत् स्थावरजङ्गमम्}
{सर्वाणि खलु भूतानि धार्यन्ते सामिकैः स्वरैः}

\twolineshloka{दीप्तायताकरुणानां मृदुमध्यमयोस्तथा}
{श्रुतीनां योऽविशेषज्ञो न स आचार्य उच्यते}

\twolineshloka{दीप्ता मन्द्रे द्वितीये च प्र चतुर्थे तथैव तु}
{अतिस्वारे तृतीये च क्रुष्टे तु करुणा श्रुतिः}

\twolineshloka{श्रुतयोऽन्या द्वितीयस्य मृदुमध्यायताः स्मृताः}
{तासामपि तु वक्ष्यामि लक्षणानि पृथक्पृथक्}

\twolineshloka{आयतात्वं भवेन् नीचे मृदुत्वं तु विपर्यये}
{स्वे स्वरे मध्यमत्वं तु तत् समीक्ष्य प्रयोजयेत्}

\twolineshloka{द्वितीये विरता या तु क्रुष्टस्य परतो भवेत्}
{दीप्तां तां तु विजानीयात् प्रथमेन मृदुः स्मृता}

\twolineshloka{अत्रैव विरता या तु चतुर्थेन प्रवर्तते}
{तथा मन्द्रे भवेद् दीप्ता साम्नश्चैव समापने}

\twolineshloka{नाविरते श्रुतिं कुर्यात् स्वरयोर्नापि चान्तरे}
{न च हस्वे च दीर्घे च न चापि घुटसञ्ज्ञके}

\twolineshloka{द्विविधा गतिः पदान्तः स्थितसन्धिः सहोष्मभिः}
{पञ्चस्वेतेषु स्थानेषु विज्ञेयं घुटसञ्ज्ञितम्}

\twolineshloka{स्वरान्तराविरतानि हस्वदीर्घघुटानि च}
{श्रुतिस्थानेष्वशेषाणि श्रुतिवत् स्वरतो भवेत्}

\twolineshloka{दीप्तामुदात्ते जानीयाद् दीप्तां च स्वरिते विदुः}
{अनुदात्ते मृदुर् ज्ञेया गान्धर्वाः श्रुतिसम्पदः}

\twolineshloka{उदात्तश्चानुदात्तश्च स्वरितप्रचिते तथा}
{निघातश्चेति विज्ञेयः स्वरभेदस्तु पञ्चधा}

\closesub

% --- Kandika 8 ---
\dnsub{॥ अष्टमी कण्डिका ॥}
\resetShloka

\twolineshloka{अत ऊर्ध्वं प्रवक्ष्यामि आर्चिकस्य स्वरत्रयम्}
{उदात्तश्चानुदात्तश्च तृतीयः स्वरितः स्वरः}

\twolineshloka{य एवोदात्त इत्युक्तः स एव स्वरितात् परः}
{प्रचयः प्रोच्यते तज्ज्ञैर्न चात्रान्यत् स्वरान्तरम्}

\twolineshloka{वर्णस्वारोऽतीतस्वारः स्वरितो द्विविधः स्मृतः}
{मात्रिको वर्ण एवं तु दीर्घस्तूच्चरितादनु}

\twolineshloka{स तु सप्तविधो ज्ञेयः स्वारः प्रत्ययदर्शनात्}
{पदेन तु स विज्ञेयो भवेद् यो यत्र यादृशः}

\twolineshloka{सप्त स्वरान् प्रयुञ्जीत दक्षिणं श्रवणं प्रति}
{आचार्यैर्विहितं शास्त्रं पुत्रशिष्यहितैषिभिः}

\twolineshloka{उच्चादुच्चतरं नास्ति नीचान् नीचतरं तथा}
{त्रैस्वर्ये स्वारसञ्ज्ञायां किंस्थानः स्वार उच्यते}

\twolineshloka{उच्चनीचस्य यन्मध्ये साधारणमिति श्रुतिः}
{तं स्वारं स्वारसञ्ज्ञायां प्रतिजानन्ति शैक्षिकाः}

\twolineshloka{उदात्ते निषादगान्धारावनुदात्त ऋषभधैवतौ}
{स्वरितप्रभवा ह्येते षड्जमध्यमपञ्चमाः}

\twolineshloka{यत्र कखपरा ऊष्मा जिह्वामूलप्रयोजना}
{तामथाज्ञापयेन्मात्रां प्रकृत्यैव तु सा कला}

\twolineshloka{जात्यः क्षेप्रोऽभिनिहितस्तैरोव्यञ्जन एव च}
{तिरोविरामः प्रश्लिष्टः पादवृत्तश्च सप्तमः}

\twolineshloka{स्वराणामहमेतेषां पृथग् वक्ष्यामि लक्षणम्}
{उद्दिष्टानां यथान्यायमुदाहरणमेव च}

\closesub
\closesection
